% The Quiet Theorems No. NN on the SAIM blog template v1.2 (pdfLaTeX + BibTeX).
% Page: default A4 reading layout. Use \documentclass[screen]{saim-blog} for a
% 160 x 240 mm screen page. `make post-screen` builds it without editing this file.
\documentclass{saim-blog}
\input{shared/metadata.tex}
\graphicspath{{figures/}{shared/brand/logos/}}
\input{shared/components/math.tex}
\input{shared/components/theorems.tex}

\hypersetup{
  pdftitle={\ProjectTitle},
  pdfauthor={\ProjectAuthorsPDF},
  pdfsubject={\ProjectSubtitle},
  pdfkeywords={\ProjectKeywords}
}

\begin{document}
\maketitle

\begin{tldr}
Familiar fact. The result, attributed to its originator with year. What it means.
\end{tldr}

% =====================================================================
\SecSetup
Probability space, objects, symbol table, numbered assumptions.

\begin{assumption}\label{as:one}
The minimal condition.
\end{assumption}

% =====================================================================
\SecResult
\begin{mainresult}
\begin{theorem}[Originator, Year]\label{thm:main}
Under \cref{as:one}, the statement.
\end{theorem}
\end{mainresult}

% =====================================================================
\SecProof
\begin{lemma}\label[lemma]{lem:one}
Auxiliary statement.
\end{lemma}
\begin{proof}
Proof of the lemma.
\end{proof}

\begin{proof}[Proof of \cref{thm:main}]
\step{First idea} Text.
\interchange{Fubini--Tonelli} Dominating function or Tonelli bound, joint measurability.
\step{Conclusion} Text.
\end{proof}

% =====================================================================
\SecInsight
\begin{enumerate}
\item \textbf{What it says.} Text.
\item \textbf{Why it is surprising.} Text.
\end{enumerate}

\begin{pullquote}
One sentence worth lifting out.
\end{pullquote}

% =====================================================================
\SecUses
One sentence on how the consequences follow.

\subsection{First consequence}
Statement and proof, using \cref{lem:one}.

% =====================================================================
\SecExample
One fully explicit computation, a table, and a TikZ figure.

% =====================================================================
\SecCheck
One sentence on the method per row.

\begin{center}
\small
\begin{tabular}{@{}lcc@{}}
\toprule
\SAIMTableHead{Case} & \SAIMTableHead{Exact} & \SAIMTableHead{Numerical}\\
\midrule
Case & value & value\\
\bottomrule
\end{tabular}
\end{center}

\paragraph{Code.} Every number in the table is printed by \texttt{code/check.py} (\cref{lst:check}; run \texttt{python3 code/check.py}), whose docstring gives the requirements.

% =====================================================================
\SecNotes
Originator history only: who, when, original title, what it contains, how it reached a wide audience.
\bibliographystyle{plainnat}
\bibliography{shared/bibliography/references}

{\let\SAIMplainauthors\ProjectAuthorsPlainDisplay
 \renewcommand{\ProjectAuthorsPlainDisplay}{\SAIMplainauthors\ProjectAuthorPostnominal}%
 \postauthorcard}

\clearpage
\appendix
\section{Code}
\lstinputlisting[style=python,basicstyle=\ttfamily\footnotesize\color{Ink},caption={Numerical check (\texttt{code/check.py}).},label={lst:check}]{code/check.py}
\end{document}
